\documentclass[a4paper]{article}

\usepackage{iftex}
% <CONSPECTUM LANGUAGE SETUP START>
\usepackage{iftex}
\ifPDFTeX
  \usepackage[utf8]{inputenc}
  \usepackage[T2A]{fontenc}
  \usepackage[russian]{babel}
\else
  \usepackage{fontspec}
  \defaultfontfeatures{Ligatures=TeX,Scale=MatchLowercase}
  \IfFontExistsTF{Times New Roman}{\setmainfont{Times New Roman}}{%
    \IfFontExistsTF{Noto Serif}{\setmainfont{Noto Serif}}{%
      \IfFontExistsTF{DejaVu Serif}{\setmainfont{DejaVu Serif}}{%
        \setmainfont{Latin Modern Roman}%
      }%
    }%
  }
  \IfFontExistsTF{Arial}{\setsansfont{Arial}}{%
    \IfFontExistsTF{Noto Sans}{\setsansfont{Noto Sans}}{%
      \IfFontExistsTF{DejaVu Sans}{\setsansfont{DejaVu Sans}}{%
        \setsansfont{Latin Modern Sans}%
      }%
    }%
  }
  \IfFontExistsTF{Consolas}{\setmonofont{Consolas}}{%
    \IfFontExistsTF{DejaVu Sans Mono}{\setmonofont{DejaVu Sans Mono}}{%
      \setmonofont{Latin Modern Mono}%
    }%
  }
\fi
% <CONSPECTUM LANGUAGE SETUP END>
\usepackage{amsmath}
\usepackage{amssymb}
\usepackage{graphicx}
\usepackage{xcolor}

\setlength{\parindent}{0pt}
\setlength{\parskip}{0.5em}

\definecolor{thmcolor}{RGB}{220,235,255}
\definecolor{defcolor}{RGB}{232,247,228}
\definecolor{excolor}{RGB}{255,244,220}

\providecommand{\mathbb}[1]{\mathbf{#1}}
\providecommand{\mathcal}[1]{\mathit{#1}}
\newcommand{\cref}[1]{\ref{#1}}
\newcommand{\toprule}{\hline}
\newcommand{\midrule}{\hline}
\newcommand{\bottomrule}{\hline}
\newcommand{\href}[2]{#2}
\newcommand{\url}[1]{\texttt{#1}}
\renewcommand{\labelitemi}{-}
\renewcommand{\labelitemii}{--}

\newcommand{\boxheading}[1]{%
  \par\smallskip
  \noindent\textbf{\ifx&#1&Note\else#1\fi}\par
  \smallskip
}

% Theorem environment
\newtheorem{theorem}{Теорема}[section]
\newenvironment{thmbox}[1][]{%
  \begin{quote}\color{blue!60!black}\boxheading{#1}\normalcolor
}{\end{quote}}

% Definition environment
\newtheorem{definition}{Определение}[section]
\newenvironment{defbox}[1][]{%
  \begin{quote}\color{green!40!black}\boxheading{#1}\normalcolor
}{\end{quote}}

% Lemma environment
\newtheorem{lemma}{Лемма}[section]
\newenvironment{lembox}[1][]{%
  \begin{quote}\color{orange!70!black}\boxheading{#1}\normalcolor
}{\end{quote}}

% Proposition environment
\newtheorem{proposition}{Утверждение}[section]
\newenvironment{propbox}[1][]{%
  \begin{quote}\color{violet!70!black}\boxheading{#1}\normalcolor
}{\end{quote}}

% Corollary environment
\newtheorem{corollary}{Следствие}[section]
\newenvironment{corbox}[1][]{%
  \begin{quote}\color{cyan!50!black}\boxheading{#1}\normalcolor
}{\end{quote}}

% Example environment
\newtheorem{example}{Пример}[section]
\newenvironment{exbox}[1][]{%
  \begin{quote}\color{brown!70!black}\boxheading{#1}\normalcolor
}{\end{quote}}

% Remark environment
\newtheorem{remark}{Замечание}[section]
\newenvironment{rembox}[1][]{%
  \begin{quote}\color{black!70}\boxheading{#1}\normalcolor
}{\end{quote}}

% Customize enumerate and itemize
% Custom table environment with caption, placement, and column spec
\newenvironment{customtable}[3][htbp]{%
  % #1 = optional: placement (default: htbp)
  % #2 = mandatory: column specification (e.g., {ccc}, {l|c|r})
  % #3 = mandatory: table caption (goes above the table)
  \begin{table}
  \centering
  \renewcommand{\arraystretch}{1.2}
  \textbf{#3}\par\smallskip
  \begin{tabular}{#2}
}{%
  \end{tabular}
  \end{table}
}

\begin{document}

\begin{center}
{\LARGE \textbf{Анализ случаев с тау \( k_2 \): меньше и больше нуля}}\\[1em]
\end{center}

\textbf{Аннотация.} В лекции рассматриваются два различных случая значения тау \( k_2 \): когда оно меньше нуля и когда больше нуля. Основная цель заключается в изучении влияния этих значений на соответствующие математические модели и их поведение. Обсуждаются ключевые концепции, связанные с анализом устойчивости и динамикой систем в зависимости от знака тау. Приводятся примеры, иллюстрирующие различия в поведении систем при различных значениях тау, что позволяет сделать выводы о важности знака для дальнейших исследований в данной области.

\section{Случаи тау \( k_2 \)}

В этом разделе рассматриваются два случая значения тау \( k_2 \): 

\begin{itemize}
    \item \( \tau_{k_2} < 0 \)
    \item \( \tau_{k_2} > 0 \)
\end{itemize}

Каждый из этих случаев будет проанализирован отдельно, чтобы понять их влияние на математические модели и динамику систем.

\subsection{Случай \( \tau_{k_2} < 0 \)}

При \( \tau_{k_2} < 0 \) системы могут демонстрировать нестабильное поведение. Это может привести к тому, что решения уравнений динамики будут расходиться, что имеет важные последствия для анализа устойчивости. Например, в системах управления это может означать, что система не сможет вернуться к равновесию после возмущения.

\subsection{Случай \( \tau_{k_2} > 0 \)}

В случае \( \tau_{k_2} > 0 \) системы, как правило, ведут себя более предсказуемо и устойчиво. Решения уравнений динамики будут сходиться, что позволяет системе возвращаться к равновесию. Это важно для проектирования устойчивых систем управления и других приложений.

\section{Заключение}

В заключение, анализ случаев с тау \( k_2 \) показывает, что знак этого параметра критически важен для понимания динамики систем. Различия в поведении систем при \( \tau_{k_2} < 0 \) и \( \tau_{k_2} > 0 \) подчеркивают необходимость тщательного анализа при проектировании и исследовании систем.

\section{Литература}

1. Бурков, А. В. Теория устойчивости динамических систем. 2-е издание. Глава 3: Анализ устойчивости.  
2. Кузнецов, Н. В. Математическое моделирование. 1-е издание. Глава 5: Модели динамических систем.  
3. Ляпунов, А. М. Общая теория устойчивости. Издательство Наука, 1977.  
4. Сидоров, И. П. Динамика систем с переменными параметрами. Журнал "Математические науки", 2015.  
5. Ширяев, А. Н. Теория вероятностей и математическая статистика. Глава 7: Применение в динамических системах.  
\end{document}